\documentclass[../Main.tex]{subfiles}

\begin{document}
\section{Preliminaries}
\subsection{Multivariate Normal Theory}
Let $\vec{X} = (X_1, \cdots, X_n)^T$ be a vector of random variables. Recall that the expectation is a vector and the variance/covariance is a matrix.
\begin{align*}
    \left[\E[\vec{X}]\right]_i &= \E[X_i] \\
    \left[\Var(\vec{X})\right]_{ij} &= \Cov(X_i, X_j)
\end{align*}

For a matrix $A \in \R^{m \times n}$ and $\vec{b} \in \R^m$,
\begin{equation*}
    E[A\vec{X} + \vec{b}] = AE[\vec{X}] + \vec{b},\quad \Cov(A\vec{X} + \vec{b}) = A\Cov(X)A^T
\end{equation*}
\begin{definition}{Multivariate normal distribution}
    A random vector $\vec{X}$ has a \underline{multivariate normal distribution} if, for any $\vec{t} \in \R^n$, the random variable $\vec{t} \cdot \vec{X}$ has a normal distribution.
\end{definition}
\begin{proposition}
    If $\vec{X}$ is a multivarite normal vector, then $A\vec{X} + \vec{b}$ is a multivariate normal for any $A \in \R^{m \times n}$ and $\vec{b} \in \R^m$.
    \label{propTransformMVN}
\end{proposition}
\begin{proof}
    Take any $\vec{t} \in \R^m$. Then:
    \begin{align*}
        \vec{t} \cdot \left(A\vec{X} + \vec{b}\right) &= (A^T \vec{t})^T \vec{X} + \vec{t} \cdot \vec{b} \\
    \end{align*}
    Then since $\vec{X}$ is a multivariate normal, $(A\vec{t}) \cdot \vec{X}$ has a normal distribuition. Adding a constant ($\vec{t} \cdot \vec{b}$) shifts the distribution but it is still normal.
\end{proof}
\begin{proposition}
    A multivariate normal distribution is fully specified by its mean $\vec{\mu}$ and covariance matrix.
    \label{propMVNParameters}
\end{proposition}
\begin{proof}
    Let $\vec{X_1}$ and $\vec{X_2}$ be multivariate normal. Let them both have mean $\vec{\mu} \in \R^n$ and covariance matrix $\Sigma \in \R^{n \times n}$. The moment generating functions are:
    \begin{align*}
        M_{\vec{X_1}}(t) &= \E\left[e^{\vec{t} \cdot \vec{X_1}}\right] \\
        &= \exp\left(\vec{t} \cdot \vec{\mu} + \frac12 \vec{t} \cdot \Sigma \vec{t}\right) \text{ by recognising normal dist.}
    \end{align*}
    and this is the same for $\vec{X_2}$, so the random variables must have the same distribution.
\end{proof}
\begin{corollary}
    Suppose $X_1, \cdots, X_n$ are independent random variables with $X_i \sim N(\mu_i, \sigma^2)$. Let $\vec{X}$ be the vector of these random variables, and $\vec{\mu}$ be the vector of the means.

    Then for an orthogonal matrix $A$, the elements of $\vec{Y} = A\vec{X}$ are independently distributed, with $Y_i \sim N((A\vec{\mu})_i, \sigma^2)$
    \label{corRotateNormals}
\end{corollary}
\begin{proof}
    The distribution of $\vec{Y}$ is determined by its mean and covariance.

    \begin{equation*}
        \E[\vec{Y}] = \E[A\vec{X}] = A\E[\vec{X}] = A\vec{\mu}.
    \end{equation*}

    \begin{align*}
        \Cov(A\vec{X}) &= A\Cov(\vec{X})A^T \\
        &= A(\sigma^2 I)A^T \\
        &= \sigma^2 AA^T \\
        &= \sigma^2 I
    \end{align*}
\end{proof}
\subsection{Distribution of \texorpdfstring{$\overline{X}$}{X-bar} and \texorpdfstring{$S_{xx}$}{Sxx}}
Suppose we have a normal distribution with unknown mean and variance. Then the samples $X_i$ are independent and identically distributed, $X_i \sim N(\mu, \sigma^2)$. Then the MLE for $\mu$ is:
\begin{equation*}
    \hat{\mu} = \overline{X} = \frac1n \sum_{i=1}^{n} X_i
\end{equation*}
and the MLE for $\sigma^2$ is:
\begin{equation*}
    \hat{\sigma}^2 = \frac{S_{xx}}{n} = \frac1n \sum_{i=1}^{n} (X_i - \overline{X})^2
\end{equation*}
\begin{theorem}
    With the above setup,
    \begin{enumerate}
        \item $\overline{X} \sim N\left(\mu, \frac{\sigma^2}{n}\right)$
        \item $\frac{S_{xx}}{\sigma^2} \sim \chi_{n-1}^2$
        \item $\overline{X}$ and $S_{xx}$ are independent.
    \end{enumerate}
    \label{thmNormalStatistics}
\end{theorem}
\begin{proof}
    Part 1 is an easy consequence of the fact that linear combinations of Gaussians are Gaussian.

    For parts 2 and 3, write:
    \begin{align*}
        \sum_{i=1}^{n} (X_i - \mu)^2 &= \sum_{i=1}^{n} \left((X_1 - \overline{X}) + (\overline{X} - \mu)\right)^2 \\
        &= \sum_{i=1}^{n} \left((X_i - \overline{X})^2 + 2(X_i - \overline{X})(\overline{X} - \mu) + (\overline{X} - \mu)^2\right) \\
        &= S_{xx} + n(\overline{X} - \mu)^2
    \end{align*}
    Let $A$ be an orthogonal matrix such that:
    \begin{align*}
        \vec{Y} &= A\left(\vec{X} - \begin{pmatrix}\mu \\ \vdots \\ \mu\end{pmatrix}\right) = \begin{pmatrix}\sqrt{n}(\overline{X} - \mu) \\ Y_2 \\ \vdots \\ Y_n\end{pmatrix} \\
        \therefore A &=
        \begin{pmatrix}
            1/\sqrt{n} & 1/\sqrt{n} & \cdots & 1/\sqrt{n} \\
             & & & \\
             & \large\tilde{A} & &
        \end{pmatrix}
    \end{align*}
    where $\tilde{A}$ is chosen such that $A$ is orthogonal.

    Then $Y_1 = \sqrt{n}(\overline{X} - \mu) \sim N(0, \sigma^2)$ and is independent of the rest of the $Y_i$ by \thmref{corRotateNormals}.
    \begin{align*}
        \norm{\vec{Y}}^2 &= \norm{A(\vec{X} - \vec{\mu})}^2 \\
        &= (\vec{X} - \vec{\mu})^T A^T A (\vec{X} - \vec{\mu}) \\
        &= (\vec{X} - \vec{\mu})^T (\vec{X} - \vec{\mu}) \\
    \end{align*}
    which means,
    \begin{equation*}
        \sum_{i=1}^{n} Y_i^2 = \sum_{i=1}^{n} (X_i - \mu)^2
    \end{equation*}
    and so we must have:
    \begin{equation*}
        \sum_{i=2}^{n} Y_i^2 = \sum_{i=1}^{n} (X_i - \mu)^2 - n(\overline{X} - \mu)^2 = S_{xx}
    \end{equation*}
    so $S_{xx}$ is independent of $Y_1 = \overline{X}$.

    Further,
    \begin{equation*}
        \frac{S_{xx}}{\sigma^2} = \sum_{i=2}^{n} \frac{Y_1^2}{\sigma^2} \sim \chi_{n-1}^2
    \end{equation*}
\end{proof}
\section{\texorpdfstring{$t$}{T}-tests}
\begin{definition}{$t$-distribution}
    If $U \sim N(0, 1)$ and $V \sim \chi_n^2$, with $U$ and $V$ independent, then
    \begin{equation*}
        T = \frac{U}{\sqrt{V / n}}
    \end{equation*} has a \underline{$t$-distribution} with $n$ degrees of freedom. We write:
    \begin{equation*}
        T \sim t_n.
    \end{equation*}
\end{definition}
\begin{example}
    Suppose that $X_i$ are independent and identically distributed with $X_i \sim N(\mu, \sigma^2)$. When $\sigma^2$ is known, we find a $100(1-\alpha)\%$ confidence interval for $\mu$ by noting that:
    \begin{equation*}
        Z = \frac{\sqrt{n}(\overline{X} - \mu)}{\sigma} \sim N(0, 1).
    \end{equation*}
    If instead $\sigma^2$ is unknown, we note:
    \begin{equation*}
        \overline{X} \sim N(\mu, \sigma^2 / n),\quad \frac{S_{xx}}{\sigma^2} \sim \chi_{n-1}^2
    \end{equation*}
    and then find a $t$-distribution:
    \begin{equation*}
        T = \frac{\sqrt{n}(\overline{X} - \mu) / \sigma}{\sqrt{S_{xx} / \sigma^2(n-1)}} \sim t_{n-1}
    \end{equation*}
    and so we can find a confidence interval as follows:
    \begin{equation*}
        \P_{\mu, \sigma^2} \left(-t_{n-1}\left(\frac{\alpha}{2}\right) \leq \frac{\sqrt{n} (\overline{X} - \mu)}{\sqrt{S_{xx} / (n-1)}} \leq t_{n-1} \left(\frac{\alpha}{2}\right)\right) = 1- \alpha
    \end{equation*}
    and so,
    \begin{equation*}
        \overline{X} \pm \sqrt{\frac{S_{xx}}{n-1}} \frac{t_{n-1}\left(\frac{\alpha}{2}\right)}{\sqrt{n}}
    \end{equation*}
    is an exact $100(1-\alpha)\%$ confidence interval for $\mu$.

    If we instead wish to test $H_0 : \mu = \mu_0$ with $\sigma^2$ unknown, then under $H_0$,
    \begin{equation*}
        T = \frac{\sqrt{n}(\overline{X} - \mu_0)}{\sqrt{S_{xx} / (n-1)}} \sim t_{n-1}
    \end{equation*}
    so reject $H_0$ when $\abs{T} \geq t_{n-1}\left(\frac{\alpha}{2}\right)$. This is a \underline{$t$-test}.
\end{example}
\begin{remarks}
    \item As $n \to \infty$, $\frac{S_{xx}}{n-1} \to \sigma^2$ and $t_{n-1}\left(\frac{\alpha}{2}\right) \to z_{\alpha / 2}$. Then the $t$-test converges to a $z$-test. Importantly, the $t$-test does not need large $n$ to be valid.
    \item The $t$-test can be seen as an example of a GLRT for:
        \begin{align*}
            H_0:& \mu = \mu_0 \\
            H_1:& \mu \neq \mu_0
        \end{align*}
        The GLR statistic is:
        \begin{equation*}
            \Lambda_{\vec{X}}(H_0, H_1) = \frac{\sup_{\mu, \sigma^2} f(\vec{x}, \mu, \sigma^2)}{\sup_{\sigma^2} f(\vec{x}, \mu_0, \sigma^2)}
        \end{equation*}
        and substituting $(\hat{\mu}, \hat{\sigma}^2) = (\overline{X}, S_{xx} / n)$ gives the numerator:
        \begin{equation*}
            \left(\frac{2\pi\sum_{i=1}^{n} (X_i - \overline{X})^2}{n}\right)^{-n / 2} \exp\left(-\frac{n}{2}\right)
        \end{equation*}
        and when $\mu = \mu_0$, the MLE for $\sigma^2$ is $\frac1n \sum_{i=1}^{n} (X_i - \mu_0)^2$. The denominator is then:
        \begin{equation*}
            \left(\frac{2\pi\sum_{i=1}^{n} (X_i - \mu_0)^2}{n}\right)^{-n / 2} \exp\left(-\frac{n}{2}\right)
        \end{equation*}
        and so the GLR statitsic is:
        \begin{align*}
            \Lambda_{\vec{X}}(H_0, H_1) &= \frac{\left(\frac{2\pi \sum_{i=1}^{n} (X_i - \overline{X})^2}{n}\right)^{-n / 2}}{\left(\frac{2\pi \sum_{i=1}^{n} (X_i - \mu_0)^2}{n}\right)^{-n / 2}} \\
            &= \left(1 + \frac{n(\overline{X} - \mu_0)^2}{\sum_{i=1}^{n} (X_i - \overline{X})^2}\right)^{-n/2}
        \end{align*}
        and so we reject $H_0$ when
        \begin{equation*}
            \frac{\sqrt{n} (\overline{X} - \mu_0)}{\sqrt{S_{xx} / (n-1)}}
        \end{equation*}
        is large. This gives another idea of where the $t$-distribution comes from.
\end{remarks}
\end{document}